\documentclass{sdxxpaper}

\papertitle{中文短文本摘要生成}
\college{计算机科学与工程学院}
\major{}
\classnumber{}
\authorname{}
\studentid{}
\reportyear{}
\reportmonth{}

\begin{document}
\makecover
\cqutBeginReportBody
\makepaperheading

\paperabstract{在此填写研究背景、方法、实验结果和主要结论。}
\paperkeywords{中文短文本摘要；深度学习；关键词}

\section{引言}
介绍研究背景、问题和本文的主要工作。

\section{数据与方法}
\subsection{数据处理}
介绍数据来源、划分和预处理方法。

\subsection{方法设计}
介绍模型结构、训练方法和推理策略。

\section{实验与结果}
\subsection{实验配置}
介绍实验环境、参数和评价指标。

\subsection{性能对比}
表~\ref{tab:results} 展示结果的排版方式，数据请按实际实验填写。
\begin{table}[htbp]
  \centering
  \caption{实验结果示例}
  \label{tab:results}
  \begin{tabular}{lccc}
    \toprule
    方法 & 指标一 & 指标二 & 指标三 \\
    \midrule
    方法 A & -- & -- & -- \\
    方法 B & -- & -- & -- \\
    \bottomrule
  \end{tabular}
\end{table}

\subsection{消融实验}
分析各组件或参数对结果的影响。

\subsection{案例分析}
结合具体案例说明模型的表现。

\section{分析与讨论}
讨论实验结果、局限性和改进方向。

\section{总结}
总结主要工作和研究结论。

% 将条目替换为实际引用的文献，并在正文使用 \cite{...}。
\begin{thebibliography}{99}
\bibitem{example} 作者. 文献标题[J]. 期刊名称, 年份, 卷(期): 页码.
\end{thebibliography}
\end{document}
